Soit $ABC$ un triangle équilatéral. Soit $A_1$, $B_1$ et $C_1$ trois points situés à l'intérieur du
triangle $ABC$ de sorte que $BA_1 = A_1C$, $CB_1 = B_1A$, $AC_1 = C_1B$ et que
\[
\widehat{BA_1C} + \widehat{CB_1A} + \widehat{AC_1B} = 480^\circ.
\]
Soit $A_2$ le point d'intersection des droites $(BC_1)$ et $(CB_1)$ ; $B_2$ le point d'intersection des
droites $(CA_1)$ et $(AC_1)$ ; $C_2$ le point d'intersection des droites $(AB_1)$ et $(BA_1)$.

Démontrer que, si le triangle $A_1B_1C_1$ est scalène, il existe deux points par lesquels passe chacun des
cercles circonscrits aux trois triangles $AA_1A_2$, $BB_1B_2$ et $CC_1C_2$.

\emph{Remarque} : un triangle est scalène lorsque ses trois côtés sont de longueurs deux à deux distinctes.
